% =====================================================================
%  MSc Dissertation — University of Glasgow, School of Computing Science
%  Mask-Augmented Multimodal Animal Behaviour Recognition
%  Compile with Tectonic (handles biber automatically):
%      tectonic --synctex --keep-logs --keep-intermediates report.tex
%  See report/README.md for the build guide.
% =====================================================================
\documentclass[12pt,a4paper]{report}

\usepackage[a4paper,margin=2.5cm]{geometry}

% --- Fonts (engine-agnostic: Times-like under pdfTeX, Times New Roman under XeTeX) ---
\usepackage{iftex}
\ifPDFTeX
  \usepackage[T1]{fontenc}
  \usepackage[utf8]{inputenc}
  \usepackage{mathptmx}
\else
  \usepackage{fontspec}
  \IfFontExistsTF{Times New Roman}{\setmainfont{Times New Roman}}{}
\fi

\usepackage{setspace}
\onehalfspacing

\usepackage{graphicx}
\usepackage{booktabs}
\usepackage{amsmath,amssymb}
\usepackage{caption}
\usepackage[dvipsnames]{xcolor}
% --- Bibliography: a self-contained thebibliography block at the end of the file.
%     No biber/bibtex needed, so Tectonic builds it with zero external tools.
%     (references.bib is kept as a data source for if you later move to biblatex.) ---

\usepackage[colorlinks=true,linkcolor=black,citecolor=black,urlcolor=NavyBlue]{hyperref}

\begin{document}

% =====================================================================
%  TITLE PAGE
% =====================================================================
\begin{titlepage}
\centering
% Drop the official crest in as report/glasgow_logo.png to render it here.
\IfFileExists{glasgow_logo.png}{\includegraphics[width=0.6\textwidth]{glasgow_logo.png}\\[2cm]}{\vspace*{2cm}}

{\LARGE\bfseries Mask-Augmented Multimodal\\[4pt] Animal Behaviour Recognition\par}
\vspace{0.6cm}
{\large Does the animal's foreground silhouette \emph{shape} help recognise behaviour,\\ where background segmentation maps did not?\par}

\vspace{2.2cm}
{\Large Kostiantyn Boiar\par}

\vspace{2.2cm}
{\large
School of Computing Science\\[2pt]
Sir Alwyn Williams Building\\[2pt]
University of Glasgow\\[2pt]
G12 8RZ\par}

\vspace{2.2cm}
{\normalsize A dissertation presented in part fulfilment of the requirements of the\\
Degree of Master of Science at the University of Glasgow\par}

\vspace{1.5cm}
{\normalsize \today\par}
\end{titlepage}

% =====================================================================
%  FRONT MATTER (roman page numbers, not listed in the ToC)
% =====================================================================
\pagenumbering{roman}

\chapter*{Abstract}
\addcontentsline{toc}{chapter}{Abstract}
% TODO: finalise once results are in.
Multimodal models recognise animal behaviour from video, and increasingly from audio and
pose. This dissertation asks a narrow, untested question: does adding the animal's
\emph{foreground silhouette}---its outline shape, tracked through a clip---improve
behaviour recognition over video(+audio) baselines, and is any benefit concentrated in the
posture- and contact-defined behaviours where shape should matter most? The motivation is a
specific published result: on the MammAlps benchmark, \emph{background/scene} segmentation
maps did not help recognition, while audio gave a small lift. Existing mask-based work uses
segmentation only to \emph{remove} the background or to crop a region of interest; none
treats the silhouette shape itself as a fused, positive input. The approach is deliberately
compute-light: every heavy encoder is frozen and its features cached once, so only a small
fusion head is trained. The contribution is the answer and its per-behaviour analysis,
rather than the model wiring.

\chapter*{Education Use Consent}
\addcontentsline{toc}{chapter}{Education Use Consent}
\noindent I hereby give my permission for this project to be shown to other University of
Glasgow students and to be distributed in an electronic form. Please note that you are under
no obligation to sign this declaration, but doing so would help future students.

\vspace{1.5cm}
\noindent\begin{tabular}{@{}p{0.5\textwidth}p{0.5\textwidth}@{}}
Name: Kostiantyn Boiar & Signature: \rule{4cm}{0.4pt}
\end{tabular}

\chapter*{Acknowledgements}
\addcontentsline{toc}{chapter}{Acknowledgements}
% TODO: thank supervisor, group, data providers, family.
I would like to thank my supervisor, \textbf{[Supervisor name]}, for their guidance and
feedback throughout this project. % Add further acknowledgements before submission.

\chapter*{AI Usage Declaration}
\addcontentsline{toc}{chapter}{AI Usage Declaration}
% TODO: complete honestly and accurately before submission (required by the School).
I acknowledge the use of AI-based tools during the preparation of this dissertation.
% Describe here exactly which tools were used and for what (e.g. literature search support,
% grammar/spelling checks, code scaffolding), and confirm that no AI tool was used to
% generate original research results, conduct experiments, or replace your own analysis.

\tableofcontents

\clearpage
\pagenumbering{arabic}

% =====================================================================
%  CHAPTER 1 — INTRODUCTION
% =====================================================================
\chapter{Introduction}
\label{ch:introduction}

\section{Motivation}
% TODO: expand. Recognising what an animal is *doing* underpins ecology and neuroscience;
% multimodal benchmarks combine video, audio and (scene) segmentation; the open question is
% whether the animal's foreground shape adds signal that the background cannot.
\textit{[Draft] Recognising what an animal is doing from camera footage underpins both
ecology and neuroscience. Recent benchmarks combine several signals---video, audio, and
reference scene-segmentation maps---yet the representation that discards the background
entirely and keeps only the animal's shape has not been tested as a positive input.}

\section{Contributions}
\begin{itemize}
  \item A reproducible, frozen-feature \textbf{video} and \textbf{video+audio} baseline on
        MammAlps Benchmark~I, scored with the authors' official evaluation script.
  \item A \textbf{foreground-silhouette shape} modality (SAM\,2 masks $\rightarrow$ shape
        descriptors) fused via a light, capacity-matched head with modality dropout.
  \item An \textbf{ablation ladder} (video $\rightarrow$ +mask $\rightarrow$ +audio
        $\rightarrow$ +audio+mask) and a \textbf{per-behaviour} analysis focused on
        posture- and contact-defined classes, plus a background-shift robustness test.
\end{itemize}

% =====================================================================
%  CHAPTER 2 — BACKGROUND AND RELATED WORK
% =====================================================================
\chapter{Background and Related Work}
\label{ch:related}

This chapter reviews the work this dissertation builds on and positions its contribution.
It covers multimodal behaviour-recognition benchmarks and the role of the background
(Sections~\ref{sec:multimodal}--\ref{sec:background}), the use of segmentation and masks in
action recognition (Section~\ref{sec:masks}), pose-based behaviour analysis
(Section~\ref{sec:pose}), and modality-fusion under missing inputs (Section~\ref{sec:fusion}),
before summarising the gap this work addresses (Section~\ref{sec:gap}).

\section{Multimodal animal behaviour recognition}
\label{sec:multimodal}
Recognising behaviour from camera footage has shifted from single-stream video classifiers
to multimodal models. MammAlps~\cite{mammalps} is a multimodal, multi-view wildlife
benchmark providing video, audio, and reference scene-segmentation maps, with a hierarchical
label space of activities and actions across several species and a VideoMAE-based
baseline~\cite{videomae}. Two of its findings frame this dissertation. First, audio gives a
small but real improvement over video alone. Second, the reference \emph{scene}
segmentation maps---a single semantic map of the background per fixed camera viewpoint,
essentially static across a clip---did \emph{not} improve over the video or video+audio
models. In other words, the benchmark already segments, but it segments the scene rather
than the animal, and that scene information does not help the model read behaviour.

\section{The role of the background}
\label{sec:background}
PanAf-FGBG~\cite{panaf} pairs each chimpanzee video with a matching empty-background video
and uses SAM\,2~\cite{sam2} masks to \emph{erase} the animal and synthesise backgrounds. Its
purpose is to expose how strongly models lean on behaviour-correlated background, and how
that reliance harms out-of-distribution generalisation across camera locations. Crucially,
the mask there is a tool for \emph{removing} the foreground: the object of study is the
background, not the animal's shape.

\section{Masks and segmentation in action recognition}
\label{sec:masks}
Mask-guided action recognition has been explored, but typically as a means of
background suppression rather than as a shape signal in its own right. MAROON~\cite{maroon},
for example, forms background-removed cutouts (via Mask R-CNN) and feeds them as an
additional stream, without treating the silhouette shape as a dedicated, fused modality and
without pose. The Segment Anything Model~2 (SAM\,2)~\cite{sam2} makes per-frame foreground
segmentation cheap and promptable---masks can be initialised from animal tracks or boxes and
propagated across a clip---which is the enabling tool for extracting a foreground silhouette
at scale.

\section{Pose-based behaviour analysis}
\label{sec:pose}
Pose-based pipelines encode the animal's skeleton rather than raw pixels. Spatial-temporal
graph networks such as ST-GCN~\cite{stgcn} model keypoint trajectories, and pose-driven
behaviour classifiers such as PoseR~\cite{poser} are lightweight and effective. Pose is
compact, but keypoints capture body \emph{configuration} only sparsely, and they weakly
capture two-body \emph{contact}, where the relevant signal is how two bodies overlap and
touch. LookAgain~\cite{lookagain} reports exactly this asymmetry: added visual context helps
some contact behaviours yet not those whose geometry pose already pins down---motivating a
dedicated shape stream for contact- and posture-defined behaviours.

\section{Missing-modality robustness and fusion}
\label{sec:fusion}
For an added modality to be practical it must be droppable at test time. ActionMAE~\cite{actionmae}
trains with modality dropout so that a multimodal model degrades gracefully when an input is
absent. This design directly informs the fusion head used here, allowing the silhouette
stream to be omitted at inference without retraining. Where a richer appearance feature is
needed, frozen self-supervised encoders such as DINOv2~\cite{dinov2} provide strong
masked-crop embeddings, held in reserve should hand-crafted shape descriptors underperform.

\section{Datasets}
\label{sec:datasets}
MammAlps~\cite{mammalps} is the primary dataset: it natively carries video, audio and
behaviour labels and supplies the published baseline to reproduce. MABe22~\cite{mabe22}
(mouse-triplet subset: clean top-view video and named contact behaviours, but no audio)
serves as a controlled second setting in which inter-animal contact geometry can be tested
directly, which MammAlps's largely single-animal clips cannot.

\section{Summary and research gap}
\label{sec:gap}
If the background does not help behaviour recognition~\cite{mammalps} and models over-rely
on it~\cite{panaf}, the obvious untested question is whether the one representation that
discards the background entirely---the animal's foreground silhouette \emph{shape}---adds
signal. Unlike scene segmentation, a per-frame foreground silhouette is dynamic: its changing
outline is posture, and the overlap of two animals' masks is contact. Treating this
silhouette as a fused, positive modality (rather than a background-removal tool or an ROI
crop) is unaddressed, and is the gap this dissertation targets.

% =====================================================================
%  CHAPTER 3 — METHODOLOGY  (skeleton)
% =====================================================================
\chapter{Methodology}
\label{ch:method}
% TODO: write up as the implementation lands.
\section{Datasets and splits}
\section{Frozen feature extraction and caching}
\section{Foreground silhouette and shape features}
\section{Fusion head and modality dropout}
\section{Training and evaluation protocol}

% =====================================================================
%  CHAPTER 4 — EVALUATION  (skeleton)
% =====================================================================
\chapter{Evaluation}
\label{ch:eval}
% TODO: results once experiments are run.
\section{Metrics}
\section{Video and video+audio baselines}
\section{Ablation ladder: video $\to$ +mask $\to$ +audio $\to$ +audio+mask}
\section{Per-behaviour analysis}
\section{Background-shift robustness}

% =====================================================================
%  CHAPTER 5 — CONCLUSION  (skeleton)
% =====================================================================
\chapter{Conclusion}
\label{ch:conclusion}
\section{Discussion}
\section{Limitations}
\section{Future work}

% =====================================================================
%  BIBLIOGRAPHY  (kept in sync with references.bib)
% =====================================================================
\addcontentsline{toc}{chapter}{Bibliography}
\begin{thebibliography}{99}

\bibitem{panaf} O.~Brookes et al. \emph{The PanAf-FGBG Dataset: Understanding the Impact of Backgrounds in Wildlife Behaviour Recognition.} CVPR, 2025. arXiv:2502.21201.

\bibitem{mammalps} V.~Gabeff, H.~Qi, B.~Flaherty, G.~S\"umb\"ul, A.~Mathis, and D.~Tuia. \emph{MammAlps: A Multi-view Video Behavior Monitoring Dataset of Wild Mammals in the Swiss Alps.} CVPR, 2025. arXiv:2503.18223.

\bibitem{lookagain} W.~Li, J.~Ke, Y.~Wang, C.~Li, and A.~Wu. \emph{Learning When to Look: On-Demand Keypoint--Video Fusion for Animal Behavior Analysis.} 2026. arXiv:2603.07279.

\bibitem{dinov2} M.~Oquab et al. \emph{DINOv2: Learning Robust Visual Features without Supervision.} Transactions on Machine Learning Research, 2024. arXiv:2304.07193.

\bibitem{poser} P.~N.~Mullen, B.~Bowlby, H.~C.~Armstrong, A.~Gray, and M.~F.~Zwart. \emph{PoseR: a deep learning toolbox for classifying animal behaviour.} Open Biology, 16(1):250322, 2026. doi:10.1098/rsob.250322.

\bibitem{sam2} N.~Ravi et al. \emph{SAM\,2: Segment Anything in Images and Videos.} ICLR, 2025. arXiv:2408.00714.

\bibitem{maroon} F.~Schindler, V.~Steinhage, S.~T.~S.~van Beeck Calkoen, and M.~Heurich. \emph{Action Detection for Wildlife Monitoring with Camera Traps Based on Segmentation with Filtering of Tracklets (SWIFT) and Mask-Guided Action Recognition (MAROON).} Applied Sciences, 14(2):514, 2024. doi:10.3390/app14020514.

\bibitem{mabe22} J.~J.~Sun et al. \emph{The MABe22 Benchmark: Learning Multi-agent and Multi-species Behavior Representations.} ICML, 2023. arXiv:2207.10553.

\bibitem{videomae} Z.~Tong, Y.~Song, J.~Wang, and L.~Wang. \emph{VideoMAE: Masked Autoencoders are Data-Efficient Learners for Self-Supervised Video Pre-Training.} NeurIPS, 2022. arXiv:2203.12602.

\bibitem{actionmae} S.~Woo et al. \emph{Towards Good Practices for Missing Modality Robust Action Recognition.} AAAI, 2023. arXiv:2211.13916.

\bibitem{stgcn} S.~Yan, Y.~Xiong, and D.~Lin. \emph{Spatial Temporal Graph Convolutional Networks for Skeleton-Based Action Recognition.} AAAI, 2018.

\end{thebibliography}

\end{document}
